\clearpage
\section*{S5 Further tables of the paper}

\begin{table}[htbp]
\centering
\small
\caption{Mean flagged fraction $\hat\alpha$ of ESF, by design and fraction of
outliers $\varepsilon$.}
\label{tab:ahat}
\begin{tabular}{lccccccc}
\toprule
Design & 0.00 & 0.05 & 0.10 & 0.20 & 0.25 & 0.30 & 0.35 \\
\midrule
D1 & 0.019 & 0.060 & 0.106 & 0.211 & 0.248 & 0.306 & 0.347 \\
D2 & 0.006 & 0.050 & 0.101 & 0.199 & 0.230 & 0.213 & 0.108 \\
D3 & 0.020 & 0.065 & 0.112 & 0.201 & 0.248 & 0.293 & 0.302 \\
D4 & 0.033 & 0.076 & 0.118 & 0.207 & 0.257 & 0.304 & 0.347 \\
D6 & 0.018 & 0.058 & 0.106 & 0.170 & 0.151 & 0.074 & 0.018 \\
D7 & 0.028 & 0.068 & 0.115 & 0.208 & 0.257 & 0.304 & 0.344 \\
D8 & 0.014 & 0.057 & 0.105 & 0.204 & 0.250 & 0.301 & 0.350 \\
D5 & 0.019 & 0.062 & 0.111 & 0.182 & 0.119 & 0.051 & 0.021 \\
\bottomrule
\end{tabular}
\end{table}

\begin{table}[htbp]
\centering
\small
\caption{Computing time per data set in seconds, the first two parts of the study pooled
(without the $\varepsilon=0.30$ cells of the first part, which the second part replaces).
TCLUST-REG ran under GNU Octave, which is slower than MATLAB, so its times overstate its cost
(by a factor of about $36$ against MATLAB Online on the data sets of Table~\ref{tab:matlab}, on different machines) and are indicative only.}
\label{tab:time}
\begin{tabular}{lcc}
\toprule
Method & Median & 90th percentile \\
\midrule
ESF & 0.08 & 1.78 \\
ESF alone & 0.07 & 1.66 \\
TCLUST-REG & 3.32 & 5.11 \\
TCLUST-REG, second trimming & 176.44 & 240.00 \\
Trimmed CWM & 4.05 & 5.42 \\
TLE & 10.42 & 27.16 \\
CTLE & 14.63 & 36.71 \\
Bisquare mixture & 3.76 & 9.13 \\
Laplace mixture & 4.46 & 9.65 \\
Contaminated normal mixture & 0.46 & 0.76 \\
Noise-component mixture & 0.02 & 0.04 \\
\bottomrule
\end{tabular}
\end{table}

\begin{table}[htbp]
\centering
\scriptsize
\caption{Plan~5: mean accuracy on clean units of ESF alone and of ESF with the
recovery step, the best competitor chosen after the run and the paired difference (standard error in
brackets), and the mean fractions of units flagged by their residuals and by the covariate screen
only ($50$ data sets per cell).}
\label{tab:plan5}
\resizebox{\textwidth}{!}{\begin{tabular}{llcccccc}
\toprule
Design & $\varepsilon$ & ESF alone & ESF & Best competitor & Difference (s.e.) & Flagged, residual & Flagged, covariates \\
\midrule
skewed errors & 0.00 & 0.923 & 0.923 & TCLUST-REG (0.15): 0.923 & -0.001 (0.001) & 0.052 & 0.000 \\
skewed errors & 0.10 & 0.917 & 0.917 & noise comp.: 0.918 & -0.001 (0.001) & 0.131 & 0.000 \\
skewed errors & 0.20 & 0.923 & 0.923 & TLE(0.25): 0.924 & -0.001 (0.001) & 0.211 & 0.000 \\
70/30, three covariates & 0.00 & 0.908 & 0.908 & bisquare: 0.911 & -0.003 (0.001) & 0.020 & 0.007 \\
70/30, three covariates & 0.10 & 0.884 & 0.906 & TCLUST-REG (0.15): 0.910 & -0.004 (0.001) & 0.109 & 0.005 \\
70/30, three covariates & 0.20 & 0.883 & 0.903 & TCLUST-REG (0.25): 0.909 & -0.006 (0.005) & 0.204 & 0.005 \\
scale growing with $|x|$ & 0.00 & 0.913 & 0.913 & TCLUST-REG (0.05): 0.914 & -0.001 (0.001) & 0.059 & 0.000 \\
scale growing with $|x|$ & 0.10 & 0.913 & 0.913 & contam. normal: 0.915 & -0.002 (0.001) & 0.130 & 0.000 \\
scale growing with $|x|$ & 0.20 & 0.909 & 0.909 & TLE(0.25): 0.910 & -0.001 (0.001) & 0.213 & 0.000 \\
\bottomrule
\end{tabular}}
\end{table}

\subsection*{S5.1 Detail from Section 4 of the paper}

\paragraph{Solver comparison.} All solvers of Section~S4 returned the same optimal values when
they finished, apart from subsamples whose optimum lies outside the coefficient bounds of the
SCIP formulation.

\paragraph{Ablation of the exact fit.} With the exact subsample fit replaced by one start of
hard alternation on the subsample, ESF alone lost up to $0.04$ with three groups or three
covariates and $0.11$ with four groups; with elemental fits, $d$ units per group fitted exactly
as in RANSAC, it lost up to $0.11$ with two unequal groups, $0.09$ with three groups or three
covariates and $0.25$ with four groups; ten starts of alternation were ahead of the exact fit by
$0.025$ in one cell of D3 and behind by $0.02$ and $0.04$ with four groups (Table~S14).

\paragraph{Constants of ESF.} The $0.975$ quantile of the covariate screen is the usual
choice for reweighting the minimum covariance determinant; with one covariate the robust
distance uses the median and the median absolute deviation.

\paragraph{Settings of the competitors added after Plan~1.} The mixture with Laplace errors
of \citet{song2014robust} was run with $20$ random starts. In the mixture of contaminated normal
regressions each group has a proportion of at least one half of good points and the rest with
an inflated scale, both estimated ($50$ starts). The noise-component mixture has a component
uniform over the range of the response enlarged by $10\%$ ($10$ starts). Sequential RANSAC
extracts the $K$ lines one at a time, each the elemental fit among $500$ random draws with the
most remaining units within $2.5$ noise scales, refitted three times on those units, which are
then removed; in D4 it is given the mean of the two noise scales. The contaminated normal and
noise-component mixtures flag a unit when its posterior probability of being bad, or of the
uniform component, exceeds one half.
In D5 and D6 the runs of TCLUST-REG first planned were also made with $5\%$ second trimming in
the covariates. FSDA offers no automatic choice of the trimming level for regression.

\paragraph{Failed fits.} CTLE failed on $36$ of the $1{,}600$ data sets of the first part used
in Section~\ref{sec:low} of the paper and on $154$ of the $2{,}400$ of the second part; counted over its
successful fits only, its worst mean accuracy up to $20\%$ of outliers is still $0.59$ in D2 and
$0.72$ in D8.

\paragraph{Variation between seeds in the comparison of Plan~1.} In D3 at $\varepsilon=0.20$,
where the largest shortfall of ESF alone against the best competitor occurs, ten runs of
ESF alone with different seeds or constants on the same $50$ data sets gave mean accuracies
between $0.885$ and $0.914$.

\paragraph{Coefficient error.} The coefficient error, the largest distance between a fitted
and a true coefficient vector after relabelling, tells the same story as the accuracy
(Table~S19). Where a method keeps its accuracy near the ceiling, its median error is between
$0.1$ and $0.4$; where it fails, the error is of the order of the distance between the lines or
larger.

\paragraph{Flagged fraction.} At low contamination the excess of the flagged fraction of the
procedure over $\varepsilon$ is largest in D4, where the noisier group overlaps the other most.
It falls well below $\varepsilon$ from $\varepsilon=0.25$ in D5 and D6 and from
$\varepsilon=0.30$ in D2. After TCLUST-REG at either conservative level, reweighting and the
recovery step, the flagged fraction differs from that of ESF by at most $0.036$ in any
cell with $\varepsilon\le0.20$.

\paragraph{Competitor effort.} With $100$ starts (Table~S21) the contaminated normal and
Laplace mixtures stayed, with three groups, at $0.86$ and $0.67$ without outliers and at about
$0.5$ and $0.61$ at $\varepsilon=0.20$; with four groups the Laplace mixture stayed below $0.73$
at every level, against about $0.95$ for ESF. The noise-component mixture with $100$
starts became worse with three groups and in D8 from $\varepsilon=0.20$, presumably because its
fits of higher likelihood absorb the outliers into a wide Gaussian component, and better with
high-leverage points. One extra start from the least-squares partition (Table~S27) repaired the
contaminated normal and noise-component mixtures without outliers (three groups: $0.89$; four
groups: $0.95$), but not with outliers, where it improved them by at most $0.04$ and in places
made them worse.

\paragraph{Cost.} The competitors run from published software took about three minutes with
second trimming in the covariates. Our implementations of the contaminated normal and
noise-component mixtures, plain ECM and EM fits, took $0.46$ and $0.02$ seconds per data set.
With $n=1000$ the running time of ESF alone grows from $0.07$ to about $0.1$ seconds,
against two to about thirty seconds for the published competitors.

\paragraph{Outliers of the designs.} Uniform outliers (D6) are redrawn while within three noise
scales of a true line, so that every outlier is a genuine one; the high-leverage outliers of D5
have responses far from both lines at their covariate values.

\paragraph{Third part of the study.} The clustered outliers of the third part sit in a tight
cloud at covariate value about $2$ and response about $20$ noise scales. With dependent
covariates the covariate screen flags two to three percentage points more of the clean units
than with independent covariates, because the joint distribution of the covariates is not
elliptical. With groups of $80\%$ and $20\%$ or $85\%$ and $15\%$ at $\varepsilon=0.30$ the first
step flags $0.26$ to $0.28$ of the units on average. With $t_3$ errors and outliers the flagged
fraction exceeds $\varepsilon$ by two to five percentage points while still flagging every
outlier (Table~S25); the flagged fraction estimates the contamination only relative to Gaussian
errors. In Plan~5 the excess of the flagged fraction over $\varepsilon$ arose mostly without
outliers, and the covariate screen flagged under $1\%$ of the units (Table~S33).

\paragraph{Confirmation of the recovery step under Plan~3.} The step changed the fit on $3.4\%$ of
the data sets, and lowered a cell mean by $0.007$ to $0.010$ only where ESF had already
broken down.

\paragraph{Repeated runs.} In the five runs with different seeds, the fraction of data sets on
which the accuracies differed by more than $0.05$ was at most $15\%$ in any one design with up to
$20\%$ of outliers, the design with clustered outliers. The losses with $B=500$ presumably have
the cause described in Section~\ref{sec:further} of the paper: the selection criterion of ESF alone can
prefer two lines on the large group to one line on each group, so a wider search finds such fits
more often, and with a small group the first step is biased.

\paragraph{Sensitivity of the recovery step to the limit on the flagged fraction.} A limit of $0.40$
instead of $\tfrac13$ gained up to $0.05$ in the two small-group designs at $\varepsilon=0.30$
and moved no other cell by more than $0.01$ (Table~S23).

\paragraph{Tone perception data.} About $60\%$ of the $85$ trials nearest the rising line lie
within $0.01$ of it and the others spread up to $0.5$.

\paragraph{The recovery step after TCLUST-REG at level $0.40$.} Apart from the uniform-noise
design, the step harms the level by at most $0.016$, in D5 at $30\%$ of outliers. With uniform
noise at $35\%$, moving each constant of the step one step up or down, and a limit of $0.25$
instead of $\tfrac13$ on the flagged fraction, did not prevent the loss; the restriction factor
$8$ recovered $0.026$ of it.

\paragraph{Fishery data, other methods.} CTLE agrees with TCLUST-REG at low levels, and the
trimmed cluster-weighted model varies with its starts at every level. With $K=3$, ESF
alone gave no stable answer for any subsample size from $m=12$ to $m=20$.

\paragraph{Fishery data, TCLUST-REG at level $0.15$.} Between the levels that treat the cheap
flows as a group and those that separate the two dearer price levels, at level $0.15$, four of
the ten starts of TCLUST-REG return two lines with the same intercept and different slopes.

\paragraph{Taxi data.} The trips treated as contamination use other tariffs (suburban
destinations, Newark airport, negotiated fares) or are invalid records (a fare, distance or
duration of zero or less, or a duration above three hours). The data and the script that draws
the twenty samples are in the repository. We did not add noise to the flat fares, since the exact
group is the structure of interest.

\paragraph{Fishery data, dependence on the seed.} Among twenty disjoint sets of five seeds,
keeping the fit of highest likelihood under the model of the recovery step returned the cheap group
and the dearer line (intercepts $2.16$ to $2.18$ and $2.78$) every time with $K=2$, and the
three price levels (intercepts $1.90$ to $1.91$, $2.49$ to $2.54$ and $2.81$ to $2.84$) every
time with $K=3$. More replicates also help: in a second set of twenty seeds, these fits were
returned by $17$, $18$ and $19$ seeds with $B=100$, $500$ and $1000$ for $K=2$, and by $14$, $18$
and $19$ for $K=3$.
